problem:  
Let $\Pi_1, \Pi_2 \subset \mathbb{R}^3$ be two intersecting planes meeting along a line $\ell$ with dihedral angle $\theta \in (0,\pi)$, and let $\mathfrak{M}(\theta)$ denote the moduli space of all properly embedded minimal annuli asymptotic to $\Pi_1 \cup \Pi_2$ near infinity. For the classical case $\theta = \pi/2$, the singly periodic Scherk surfaces represent members of $\mathfrak{M}(\pi/2)$, forming a one–parameter family corresponding to separation distance or scaling.  

Consider now the following general problem:

**Question:**  
For general $\theta \in (0,\pi)$, does there exist a *nontrivial continuous family* of minimal annuli in $\mathfrak{M}(\theta)$ parametrized by a twist angle $\phi \in S^1$ such that:
1. each member $M_{\theta,\phi} \in \mathfrak{M}(\theta)$ is properly embedded and asymptotic to $\Pi_1 \cup R_\phi(\Pi_2)$, where $R_\phi$ denotes rotation about $\ell$ by angle $\phi$;
2. $\phi \mapsto M_{\theta,\phi}$ is smooth and nonconstant (so that distinct $\phi$ yield distinct minimal annuli up to isometry fixing $\ell$);
3. the corresponding Jacobi operator on each $M_{\theta,\phi}$ possesses a nontrivial kernel generated by infinitesimal twisting of the asymptotic boundary planes?

Moreover, if such a family $\{M_{\theta,\phi}\}$ exists, can these twisted annuli serve as local models in desingularization constructions within $S^3$, yielding new one–parameter moduli families of compact embedded minimal surfaces obtained by gluing along transverse intersection curves?

Why is it a "good" problem:  
This version precisely targets the key analytic and geometric issue omitted from earlier formulations—whether minimal surfaces interpolating between intersecting planes admit an additional "twist" degree of freedom.  

1. **Profound Insight:**  
The problem draws attention to the possible rigidity or flexibility of planar–end minimal surfaces. Understanding whether twisting preserves minimality probes directly the spectral and deformation theory of minimal surfaces with multi–planar asymptotics. A positive result would reveal a new source of moduli extending beyond known Scherk–type families.  

2. **Bridge Across Fields:**  
It connects geometric PDE theory (existence of minimal surfaces with prescribed asymptotics) and analysis of elliptic operators (Jacobi field classification) with global geometry in compact manifolds (using these local models for gluing in $S^3$). Resolving the existence or absence of such twisted families would impact both the local deformation theory in $\mathbb{R}^3$ and the global gluing constructions in $S^3$.  

3. **Extreme Simplicity:**  
The statement is concise: Do minimal surfaces asymptotic to two planes meeting along a line support a continuous twisting modulation? Yet it conceals deep analytical subtleties involving boundary behavior, Fredholm properties, and uniqueness.  

4. **Potential Impact:**  
If such twisted annuli exist, they would immediately enrich the local building blocks of all known desingularization frameworks, possibly leading to new continuous moduli of compact minimal surfaces in $S^3$. Conversely, proving nonexistence would establish a rigidity theorem characterizing the Scherk families as the unique local models. Either outcome would substantially advance understanding of minimal surface moduli and gluing mechanisms.  

Hence, this problem is both sharply defined and structurally deep: it isolates a pivotal mechanism—existence of twisted planar–end minimal annuli—on which future global constructions in $S^3$ would critically depend.